\documentclass[aspectratio=169,11pt]{beamer}

\usetheme{default}
\usefonttheme{professionalfonts}
\usepackage[T1]{fontenc}
\usepackage[default]{sourcesanspro}
\usepackage{amsmath,amssymb}
\usepackage{booktabs,tabularx,array,colortbl}
\usepackage{pgfplots}
\pgfplotsset{compat=1.17}

\definecolor{Ink}{HTML}{18212B}
\definecolor{Muted}{HTML}{5D6875}
\definecolor{Terra}{HTML}{9C3E2F}
\definecolor{Blue}{HTML}{2F6690}
\definecolor{Pale}{HTML}{EEF2F4}
\definecolor{Green}{HTML}{2C6E49}

\setbeamercolor{normal text}{fg=Ink,bg=white}
\setbeamercolor{frametitle}{fg=Ink,bg=white}
\setbeamercolor{alerted text}{fg=Terra}
\setbeamercolor{structure}{fg=Blue}
\setbeamercolor{footline}{fg=Muted,bg=white}
\setbeamertemplate{navigation symbols}{}
\setbeamertemplate{itemize item}{\raise0.15ex\hbox{\scriptsize\textbullet}}
\setbeamertemplate{itemize subitem}{\raise0.15ex\hbox{\scriptsize--}}
\setbeamerfont{frametitle}{size=\Large,series=\bfseries}
\setbeamerfont{title}{size=\huge,series=\bfseries}
\setbeamerfont{subtitle}{size=\large}
\setbeamertemplate{frametitle}{%
  \vspace{0.28em}\insertframetitle\par
  \vspace{0.12em}{\color{Terra}\rule{\textwidth}{0.8pt}}
}
\setbeamertemplate{footline}{%
  \leavevmode\hbox{\begin{beamercolorbox}[wd=\paperwidth,ht=2.5ex,dp=0.8ex,rightskip=1.2em]{footline}%
    \hfill\scriptsize\insertframenumber\,/\,\inserttotalframenumber
  \end{beamercolorbox}}
}

\newcommand{\source}[1]{\vfill{\color{Muted}\tiny Source: #1}}
\newcommand{\good}[1]{{\color{Green}\textbf{#1}}}
\newcommand{\bad}[1]{{\color{Terra}\textbf{#1}}}
\newcommand{\tightitems}{\setlength{\itemsep}{0.55em}}
\newcolumntype{Y}{>{\raggedright\arraybackslash}X}

\title{Calibration of a Lifecycle Model of Housing and Fertility}
\subtitle{External inputs, internal moments, and model fit}
\author{Tommaso De Santo}
\date{July 2026}

\begin{document}

% 1
\begin{frame}[plain]
  \titlepage
\end{frame}

% 2
\begin{frame}{Calibration objective}
  \large
  The model must jointly reproduce four lifecycle choices:
  \begin{itemize}\tightitems
    \item completed family size;
    \item tenure and occupied housing;
    \item liquid saving and entry wealth;
    \item late-life estates and housing retention.
  \end{itemize}

  \vspace{0.5em}
  The central mechanism is that children raise housing needs, while access to owner housing depends on wealth and the down payment.

  \vspace{0.5em}
  Calibration therefore uses variation across age, tenure, family status, housing size, and wealth.
  \source{Circulated calibration note, Section 1.}
\end{frame}

% 3
\begin{frame}{Empirical discipline}
  \begin{columns}[T,onlytextwidth]
    \column{0.31\textwidth}
      \large\textbf{CPS}
      \begin{itemize}\tightitems
        \item Completed fertility
        \item Completed childlessness
      \end{itemize}
    \column{0.31\textwidth}
      \large\textbf{ACS}
      \begin{itemize}\tightitems
        \item Ownership by age and family status
        \item Rooms by tenure and family size
        \item Aggregate occupied rooms
      \end{itemize}
    \column{0.31\textwidth}
      \large\textbf{PSID}
      \begin{itemize}\tightitems
        \item Entrant wealth
        \item Young liquid wealth
        \item Late-life estates and nonhousing wealth
      \end{itemize}
  \end{columns}

  \vspace{1.0em}
  Each dataset disciplines a different margin, but all parameters are estimated jointly.
  \source{Circulated calibration note, Sections 1.1--1.3.}
\end{frame}

% 4
\begin{frame}{Lifecycle timing}
  \large
  A model period is four years.

  \vspace{1.0em}
  \begin{center}
  \begin{tabular}{c@{\hspace{1.3em}}c@{\hspace{1.3em}}c@{\hspace{1.3em}}c}
    \textbf{18} & \textbf{fertile years} & \textbf{66} & \textbf{85}\\
    entry & family-size choice & retirement & maximum age\\
  \end{tabular}
  \end{center}

  \vspace{1.2em}
  \begin{itemize}\tightitems
    \item Survival is certain before retirement and follows the 2023 Social Security life table thereafter.
    \item Dependent children remain at home for 18 years on average: maturity probability $2/9$ per period.
    \item Risk aversion is fixed at $\sigma=2$.
  \end{itemize}
  \source{Circulated calibration note, ``Demographics and preferences.''}
\end{frame}

% 5
\begin{frame}{Preferences and family needs}
  \large
  Children affect both the utility from family size and the resources required by the household.

  \vspace{0.7em}
  Consumption and housing requirements are
  \[
    \bar c(n_s)=\bar c_0+\bar c_n n_s,
    \qquad
    \bar h(n_s)=\bar h_0+\bar h_{jump}\mathbf 1\{n_s\geq 1\}+\bar h_n n_s.
  \]

  \begin{itemize}\tightitems
    \item $\psi_{child}$ governs the direct utility value of children.
    \item $\kappa_F$ governs dispersion in completed-family-size choices.
    \item $\bar h_{jump}$ isolates the first-child increase in space needs.
  \end{itemize}
  \source{Circulated calibration note, Tables 2--3 and ``Consumption needs and family size.''}
\end{frame}

% 6
\begin{frame}{Housing products}
  \begin{columns}[T,onlytextwidth]
    \column{0.47\textwidth}
      \large\textbf{Renting}
      \[
        h^R\in[0,6].
      \]
      Renters choose housing continuously up to six rooms, the 90th percentile of the matched ACS rental distribution.
    \column{0.47\textwidth}
      \large\textbf{Owning}
      \[
        H\in\{2,4,6,8,10\}.
      \]
      Owners choose from five discrete house sizes that summarize the observed room distribution.
  \end{columns}

  \vspace{1.0em}
  The tenure-specific menus make family-sized housing lumpy and link fertility to tenure choice.
  \source{Circulated calibration note, ``Housing products and supply.''}
\end{frame}

% 7
\begin{frame}{Housing finance}
  \large
  Mortgage debt cannot exceed $\phi=0.80$ of house value, so purchase requires
  \[
    (1-\phi)PH=0.20PH
  \]
  in equity.

  \vspace{0.4em}
  \begin{columns}[T,onlytextwidth]
    \column{0.47\textwidth}
      \begin{itemize}\tightitems
        \item Real liquid return: 2.0\%
        \item Housing depreciation: 1.1\%
        \item Property tax: 1.0\%
      \end{itemize}
    \column{0.47\textwidth}
      \begin{itemize}\tightitems
        \item Selling cost: 6\%
        \item Labor-income tax: 17.9\%
        \item No separate payment-to-income limit
      \end{itemize}
  \end{columns}
  \source{Circulated calibration note, ``Housing finance and transaction costs.''}
\end{frame}

% 8
\begin{frame}{Income process}
  \large
  Earnings combine a deterministic age profile with persistent idiosyncratic risk:
  \par\medskip
  \[
    \log z_{t+1}=\rho_z\log z_t+\varepsilon_{t+1},
    \qquad
    \rho_z=0.960,
    \qquad
    \sigma_z=0.0645.
  \]

  The annual process is sampled at the four-year model frequency and approximated by a five-state Rouwenhorst chain.

  \vspace{0.8em}
  The persistence is consistent with the upper end of PSID estimates. The innovation variance is deliberately below the empirical range and the quantitative results are conditional on it.
  \source{Circulated calibration note, ``Income process.''}
\end{frame}

% 9
\begin{frame}{Housing supply}
  \large
  Aggregate housing services are supplied elastically, with elasticity
  \[
    \eta=1.75,
  \]
  based on the population-weighted metropolitan estimate in Saiz (2010).

  \vspace{0.8em}
  The scale of housing supply, $\bar H$, is estimated internally so that the model matches aggregate occupied rooms.

  \vspace{0.8em}
  The calibration therefore separates:
  \begin{itemize}\tightitems
    \item an externally assigned supply elasticity; and
    \item an internally estimated supply level.
  \end{itemize}
  \source{Circulated calibration note, ``Housing products and supply.''}
\end{frame}

% 10
\begin{frame}{Entry wealth and bequests}
  \large
  Households enter at age 18 with a draw from the PSID net-worth-to-income distribution of childless renters aged 18--24.

  \vspace{0.8em}
  The benchmark imposes
  \[
    \theta_n=0,
  \]
  so the bequest motive does not vary directly with the number of children.

  \vspace{0.8em}
  Children can still affect estates indirectly through housing needs, tenure, and lifecycle saving.
  \source{Circulated calibration note, ``Entry wealth and bequests.''}
\end{frame}

% 11
\begin{frame}{Parameters set outside estimation}
  \small
  \centering
  \begin{tabularx}{\textwidth}{YcY}
    \toprule
    Parameter & Value & Basis\\
    \midrule
    Period; entry; retirement; max age & 4; 18; 66; 85 & Model timing\\
    Child duration; risk aversion & 18 years; $\sigma=2$ & Maintained restrictions\\
    Return; depreciation; property tax & 2.0\%; 1.1\%; 1.0\% & External evidence / maintained tax\\
    Financed share; selling cost & 0.80; 6\% & Housing-finance evidence\\
    Income persistence; innovation s.d. & 0.960; 0.0645 & PSID literature / provisional\\
    Owner menu; renter maximum & $\{2,4,6,8,10\}$; 6 & Matched ACS distributions\\
    Supply elasticity & 1.75 & Saiz (2010)\\
    Entrant wealth; $\theta_n$ & PSID distribution; 0 & Data / maintained restriction\\
    \bottomrule
  \end{tabularx}
  \source{Circulated calibration note, Table 1.}
\end{frame}

% 12
\begin{frame}{Internal estimation}
  \large
  Fourteen parameters are estimated jointly using fifteen moments.

  \vspace{0.9em}
  \begin{columns}[T,onlytextwidth]
    \column{0.23\textwidth}
      \textbf{Saving}\\[0.35em]
      $\beta_{annual}$
    \column{0.23\textwidth}
      \textbf{Housing}\\[0.35em]
      $\alpha_c,\bar h_0,$\\
      $\bar h_{jump},\bar h_n$
    \column{0.23\textwidth}
      \textbf{Family size}\\[0.35em]
      $\bar c_0,\bar c_n,$\\
      $\psi_{child},\kappa_F$
    \column{0.23\textwidth}
      \textbf{Tenure / estates}\\[0.35em]
      $\chi_O,\bar H,$\\
      $\theta_0,\theta_1,\kappa_T$
  \end{columns}

  \vspace{1.3em}
  The blocks organize the economics; the estimator chooses all parameters simultaneously because most moments load on several margins.
  \source{Circulated calibration note, Table 2 and Section 1.2.}
\end{frame}

% 13
\begin{frame}{Saving block}
  \large
  The annual discount factor is estimated at
  \[
    \beta_{annual}=0.9912.
  \]

  It is disciplined jointly by wealth at three lifecycle stages:
  \begin{itemize}\tightitems
    \item liquid wealth among young childless renters;
    \item the share of older households with substantial nonhousing wealth;
    \item the late-life median estate.
  \end{itemize}

  The estimate is a joint calibration outcome, not a direct measure of household patience.
  \source{Circulated calibration note, ``Saving.''}
\end{frame}

% 14
\begin{frame}{Housing-demand block}
  \begin{columns}[T,onlytextwidth]
    \column{0.48\textwidth}
      \large\textbf{Estimates}
      \[
      \begin{aligned}
        \bar h_0&=0.3915,\\
        \bar h_{jump}&=1.6021,\\
        \bar h_n&=0.1644.
      \end{aligned}
      \]
    \column{0.48\textwidth}
      \large\textbf{Empirical discipline}
      \begin{itemize}\tightitems
        \item Childless-renter mean rooms
        \item First-child increase in rooms
        \item Rooms gap between larger and smaller families
      \end{itemize}
  \end{columns}

  \vspace{0.7em}
  The first child creates the largest discrete increase in required space.
  \source{Circulated calibration note, ``Housing demand, tenure, and supply.''}
\end{frame}

% 15
\begin{frame}{Tenure and supply block}
  \begin{columns}[T,onlytextwidth]
    \column{0.47\textwidth}
      \large\textbf{Tenure demand}
      \[
        \chi_O=1.1133,
        \qquad
        \kappa_T=0.0100.
      \]
      Ownership rates by age and family status discipline the owner-service premium and tenure-choice dispersion.
    \column{0.47\textwidth}
      \large\textbf{Housing supply}
      \[
        \bar H=8.6456.
      \]
      Aggregate occupied rooms disciplines the supply scale; the owner-renter rooms gap helps discipline housing demand.
  \end{columns}

  \vspace{0.8em}
  Tenure and housing-size moments identify a common block because ownership changes both access and the menu of available space.
  \source{Circulated calibration note, ``Housing demand, tenure, and supply.''}
\end{frame}

% 16
\begin{frame}{Family-size block}
  \small
  \begin{columns}[T,onlytextwidth]
    \column{0.48\textwidth}
      \large\textbf{Needs}
      \[
      \begin{aligned}
        \bar c_0&=1.2596,\\
        \bar c_n&=0.3929.
      \end{aligned}
      \]
      Young childless wealth is particularly informative about the baseline requirement; family-composition differences discipline the incremental cost.
    \column{0.48\textwidth}
      \large\textbf{Preferences}
      \[
      \begin{aligned}
        \psi_{child}&=0.1979,\\
        \kappa_F&=2.0434.
      \end{aligned}
      \]
      Completed fertility and childlessness discipline the utility value of children and dispersion in family-size choices.
  \end{columns}

  \vspace{0.6em}
  Costs and tastes are estimated jointly because the same fertility outcomes respond to both.
  \source{Circulated calibration note, ``Consumption needs and family size.''}
\end{frame}

% 17
\begin{frame}{Bequest block}
  \large
  The strength and wealth dependence of the bequest motive are
  \[
    \theta_0=0.3118,
    \qquad
    \theta_1=0.3973.
  \]

  \begin{itemize}\tightitems
    \item The median late-life estate is particularly informative about $\theta_0$.
    \item Older nonhousing wealth and old-age ownership add discipline for $\theta_1$.
    \item The same moments also respond to saving and tenure, so the bequest block is estimated jointly with them.
  \end{itemize}
  \source{Circulated calibration note, ``Bequests.''}
\end{frame}

% 18
\begin{frame}{Family size and family housing}
  \small
  \centering
  \begin{tabularx}{\textwidth}{Yrrr}
    \toprule
    Moment & Data & Model & Gap\\
    \midrule
    Completed fertility & 1.918 & 2.034 & +0.116\\
    Childlessness rate & 0.188 & 0.189 & +0.001\\
    Family ownership gap & 0.168 & 0.219 & +0.052\\
    First-child housing increment, rooms & 0.664 & 0.681 & +0.016\\
    Rooms gap, 3+ versus 1--2 children & 0.368 & 0.607 & +0.240\\
    \bottomrule
  \end{tabularx}

  \vspace{1.0em}
  The model closely matches childlessness and the first-child housing response. It overstates the separation in ownership and rooms between larger and smaller families.
  \source{Circulated calibration note, Table 3.}
\end{frame}

% 19
\begin{frame}{Tenure, housing, and wealth}
  \scriptsize
  \centering
  \begin{tabularx}{\textwidth}{Yrr@{\hspace{1.5em}}Yrr}
    \toprule
    Moment & Data & Model & Moment & Data & Model\\
    \midrule
    Ownership, ages 30--55 & .575 & .658 & Ownership, ages 25--34 & .341 & .274\\
    Childless-renter rooms & 3.805 & 3.963 & Ownership, ages 65--75 & .764 & .954\\
    Childless-owner share, 6+ rooms & .596 & .760 & Mean occupied rooms & 5.780 & 5.673\\
    Owner-renter rooms gap & 2.419 & 2.323 & Young liquid wealth / income & .179 & .328\\
    Median estate / income & 6.501 & 6.431 & Older nonhousing wealth share & .608 & .610\\
    \bottomrule
  \end{tabularx}

  \vspace{0.9em}
  \normalsize
  Housing quantities and late-life wealth levels are close. The largest misses are excessive old-age ownership, too many childless owners in large homes, and excess liquid wealth among young renters.
  \source{Circulated calibration note, Table 3 and discussion of fit.}
\end{frame}

% 20
\begin{frame}{Summary}
  \large
  \begin{itemize}\tightitems
    \item The calibration combines externally assigned timing, finance, income, and supply parameters with 14 internally estimated parameters.
    \item Fifteen moments jointly discipline fertility, housing needs, tenure, saving, and bequests.
    \item The model matches childlessness, the first-child housing response, aggregate rooms, and targeted late-life wealth levels.
    \item The main tensions are ownership by age, large-home selection, young liquid wealth, and the thin upper tail of estates.
  \end{itemize}

  \vspace{0.7em}
  The current calibration is designed for quantitative exercises on family housing needs, tenure, and the down-payment constraint; conclusions that depend on late-life portfolio composition remain outside its scope.
  \source{Circulated calibration note, Section 1.3.}
\end{frame}

\end{document}
